\section{Composition of heterogeneous certificates}
\label{app:composition}

We prove Proposition~\ref{prop:composition}.  Let
$F:\Rep_0\to\Rep_1$ and $G:\Rep_1\to\Rep_2$.  Write $\omega_0$ for an
admissible source state and $B\in\mathcal D_2$ for a target observable for
which both pullbacks are defined.  Insert the intermediate pairing:
\begin{align*}
&\left|\langle(G\circ F)_*\omega_0,B\rangle
 -\langle\omega_0,F^*G^*B\rangle\right|\\
&\quad\leq
 \left|\langle G_*F_*\omega_0,B\rangle
        -\langle F_*\omega_0,G^*B\rangle\right|\\
&\qquad+
 \left|\langle F_*\omega_0,G^*B\rangle
        -\langle\omega_0,F^*G^*B\rangle\right|\\
&\quad\leq \varepsilon_G(B;i)+\varepsilon_F(G^*B;i).
\end{align*}
This proves Equation~\eqref{eq:composition}.  If
$\varepsilon_F(A;i)\leq e_F(i)p_1(A)$ and
$p_1(G^*B)\leq L_G(B)p_2(B)$, the second term is at most
$e_F(i)L_G(B)p_2(B)$.

For a third arrow $H:\Rep_2\to\Rep_3$, scalar channel notation gives
\begin{align*}
 \varepsilon_{H\circ(G\circ F)}
 &\leq \varepsilon_H+L_H(\varepsilon_G+L_G\varepsilon_F),\\
 \varepsilon_{(H\circ G)\circ F}
 &\leq (\varepsilon_H+L_H\varepsilon_G)+L_HL_G\varepsilon_F.
\end{align*}
The expressions coincide.  Observable-dependent moduli obey the same identity
with each pullback written explicitly.  Identity arrows have zero modulus and
unit Lipschitz factor.

An exact channel in the composite must be exact on both legs.  A uniformity
domain is the intersection of the two domains after translating the target
indices and supports through $G$.  Suppose a prerequisite $P$ is unavailable
after $F$.  The equation above can still control expectations of observables
in the pullback domain, but it cannot establish $P$ unless either (i) $P$ is
not needed for the final claim, or (ii) a recovery theorem takes available
source/intermediate data to $P$ on the required domain.  This is the formal
meaning of the rule that a forgotten certificate does not reappear by
composition alone.

Associativity is therefore established for the typed maps and quantitative
moduli.  Recovery obligations compose as a dependency graph: the output type
of each recovery theorem must match the prerequisite type of the next
inference.  We do not replace that graph by a Boolean meet when the theorem
changes the bearer of the property.
